\documentclass[10pt,a4paper]{amsart}
\usepackage[margin=19mm]{geometry}
\usepackage{amsmath,amssymb,mathtools}
\usepackage[scr=rsfs,cal=euler]{mathalfa}
\usepackage{xfrac}
\usepackage[unicode,hidelinks,bookmarks=false]{hyperref}
\newcommand{\ie}{i.e.\ }
\newcommand{\cf}{cf.\ }
\newcommand{\resp}{respectively\ }
\newcommand{\Subsection}{\subsection}
\providecommand{\nobreakdash}{\nobreak\hbox{-}}
\setlength{\emergencystretch}{2em}
\multlinegap0pt
\usepackage{amssymb}
\usepackage{mathtools}
\usepackage[shortlabels]{enumitem}
\usepackage{tikz}
\usepackage{xcolor}
\newtheorem{proposition}{Proposition}
\newcommand{\D}{\mathbb D}
\title{On Cluster Deep Loci in Double Bruhat Cells}
\author{Quan Hanwen}
\date{Source: arXiv:2608.27870v1 (CC BY 4.0)}
\begin{document}
\maketitle

\noindent\emph{Source and licence.} On Cluster Deep Loci in Double Bruhat Cells. Version arXiv:2608.27870v1 (CC BY 4.0), distributed by arXiv under the Creative Commons Attribution 4.0 International licence, http://creativecommons.org/licenses/by/4.0/, as recorded in the arXiv licence metadata for this exact version and shown on the abstract page https://arxiv.org/abs/2608.27870v1. Full licence notice: \texttt{licenses/arxiv-2608.27870-CC-BY-4.0.txt}.\par\smallskip

\noindent\textit{Editorial scope.} Counted unit: the article's proposition proposition (source0.tex lines 2131--2146). The article's complete original proof of it is delivered with it (source0.tex lines 2148--2201, one \texttt{proof} environment of 54 lines). No mathematics is written, repaired or re-derived: the statement and the proof are the article's own lines, in the article's own notation, and no retained line is substituted or re-flowed.  No further source-local context is needed: every cross-reference of every retained line resolves inside the two delivered spans.  Nothing was omitted from any retained span, so the package carries no omission record.  No retained line contains a citation, so the wrapper carries no bibliography and no bibliography entry.  No retained line contains a figure environment, an image-inclusion command or drawing code, so no image is delivered and the package declares no asset dependency.  Editorial formatting only, in addition: an AMS \texttt{amsart} preamble that restates the article's own declarations and macros used by the retained lines, namely the environments \texttt{proposition} on separate counters, together with the standard packages those lines need.  The retained environments print in this document as Proposition 1 in order of appearance, whereas the article prints the same items with the numbering of its own full document.  The complete native source of the article as distributed in the arXiv e-print for this exact version is bundled unchanged as \texttt{sources/208745/source0.tex}.
\begin{proposition}
   For a pair of double Weyl group elements $(u,v)$ and $(u',v')$ with
\[
(u',v') < (u,v),\qquad \ell(u',v')=\ell(u,v)-1
\]
one has
\[
G^{u',v'}
\cap
\overline{\D(G^{u,v})}
\subseteq
\D_{\mathrm{ext}}^{(u',v')\subset(u,v)}.
\]
where $\D_{\mathrm{ext}}^{(u',v')\subset(u,v)}$ is the complement of all extendable coordinate tori of $(u',v')$ in $G^{u',v'}$.

\end{proposition}

\begin{proof}
Suppose that
\[
x\in
G^{u',v'}
\setminus
\D_{\mathrm{ext}}^{(u',v')\subset(u,v)}
\cap
\overline{\D(G^{u,v})}
\]
Then there exists an extendable word
\(
\mathbf j\in R_{\mathrm{ext}}^{(u,v)}(u',v')
\)
such that
\(
x\in T_{\mathbf j}.
\)


Choose an extension of $\mathbf j$
\[
\mathbf i=(i_1,\ldots,i_{m+1})\in R(u,v)
\]
with
\[
\mathbf j=(i_1,\ldots\hat{i_k},\ldots,i_{m+1}).
\]
 Allow the
$k$-th factorization parameter to vanish and consider
\[
\widetilde{x}_{\mathbf i,\mathbf j}:H\times(\mathbb C^*)^{k-1}\times\mathbb C\times(\mathbb C^*)^{m+1-k}\longrightarrow G^{u',v'}\sqcup G^{u,v},
\]
defined by
\[
\widetilde{x}_{\mathbf i,\mathbf j}
(h;t_1,\ldots,t_{m+1})
=
h x_{i_1}(t_1)\cdots x_{i_{m+1}}(t_{m+1}).
\]

For $t_k\neq0$, this is the factorization chart $T_{\mathbf i}$ in the higher-dimensional cell. For $t_k=0$, the factor $x_{i_k}(0)$ is the identity, and the map restricts to the factorization chart $T_{\mathbf j}$ in the boundary cell. The map is injective: its restrictions to $t_k=0$ and $t_k\neq0$ are injective factorization maps, and their images lie in two
disjoint double Bruhat cells. Moreover, the source and the target
\[
G^{u',v'}\sqcup G^{u,v}
\]
are both real manifolds of the same dimension. Therefore, by invariance of domain,
$\widetilde{x}_{\mathbf i,\mathbf j}$ is an open embedding onto its image.
In particular,
\[
T_{\mathbf j}\sqcup T_{\mathbf i}
\]
is an open neighbourhood of $x$, which is disjoint from $\D(G^{u,v})$. This contradicts the assumption that $x$ lie in the closure of $\D(G^{u,v})$.
\end{proof}

\end{document}
